{\bf An overview of Sections \ref{sec:Goodwillie}-\ref{sec:Toda}}
Fix an integer $n \ge 3$.  Recall that $\MO \langle 4n \rangle$ is, by definition, the Thom spectrum \cite{BMMS,ABGHR} of the composite spectrum map
\begin{equation} \label{eq:ThomMap}
\tau_{\ge 4n} ko \to \tau_{\ge 1}ko \to \Sigma \mathrm{gl}_1(\mathbb{S}).
\end{equation}
Here, the spectrum map $\tau_{\ge 1}ko \to \Sigma \mathrm{gl}_1(\mathbb{S})$ is the infinite delooping of the real $J$ homomorphism $\mathrm{BO} \to \mathrm{BGL}_1(\mathbb{S})$, which exists because the one-point compactification functor $\mathrm{Vect}_{\mathbb{R}} \to \mathrm{Top}_*$ is symmetric monoidal.


Our first aim in this paper is to prove \Cref{thm:intro-main}, or equivalently to understand the unit map
$$\pi_{8n-1} \bS \to \pi_{8n-1} \MO \langle 4n \rangle.$$
To begin to do so, we fix some notation and recall more precisely how a Thom spectrum such as $\MO \langle 4n \rangle$ is defined.

\begin{dfn}
Taking $\Omega^{\infty+1}$ of the sequence (\ref{eq:ThomMap}) gives maps of infinite loop spaces
$$\Omega \Omega^{\infty} \tau_{\ge 4n} ko \to \mathrm{O} \to \GL_1(\mathbb{S}).$$
We use the notation $\Of$ to denote the infinite loop space $\Omega \Omega^{\infty} \tau_{\ge 4n} ko$.
The infinite loop map
$$\Of \to \GL_1(\mathbb{S})$$
gives rise, by the universal property of $\GL_1(\mathbb{S})$ \cite[Theorem 5.1]{ABGHR}, to a map of $\E_{\infty}$-rings
$$J_+:\spOf \to \bS.$$
Here, $\spOf$ is a spherical group ring, with underlying spectrum the suspension spectrum of the pointed space $\Of_+$.
Contracting $\Of$ to a point gives rise to a second $\mathbb{E}_\infty$-ring map
$$\epsilon:\spOf \to \bS,$$
the \emph{augmentation map.}  We use $J:\sOf \to \bS$ to refer to the composite 
$$\sOf = \fib(\epsilon) \longrightarrow \spOf \stackrel{J_+}{\longrightarrow} \bS,$$
which is a map of non-unital $\mathbb{E}_\infty$-rings.
\end{dfn}

\begin{cnstr} \label{cnstr:barconstruction}
By Definition 4.1 of \cite{ABGHR}, the spectrum $\MO \langle 4n \rangle$ can be presented as the geometric realization of the two-sided bar construction
$$\MO \langle 4n \rangle \simeq \left|\sBar(\bS,\spOf, \bS)_{\bullet} \right|,$$
which computes the relative tensor product $\bS \otimes_{\spOf} \bS$.
Here, the action of $\spOf$ on the left copy of $\bS$ is via $\epsilon$, and the action on the right copy of $\bS$ is via $J_+$.
\end{cnstr}

One may view our work in the first half of the paper as a computation of the map $\pi_{8n-1} \Ss \to \pi_{8n-1} \MO \langle 4n \rangle$ via the associated bar spectral sequence. We will see that the only possible differentials affecting $\pi_{8n-1} \MO \langle 4n \rangle$ are $d_1$-differentials and a single $d_2$-differential, so what we need to do may be summarized as follows:
\begin{enumerate}
    \item Compute the $\mathrm{E}_1$-page in the relevant range.
    \item Compute the relevant $d_1$-differentials.
    \item Compute the single relevant $d_2$-differential.
\end{enumerate}
Later in this section, we study the homotopy of $\sOf$ using the Goodwillie tower in augmented $\mathbb{E}_\infty$-ring spectra.  This is enough to resolve (1) and most of (2) above.

The key to proving \Cref{thm:intro-main} is to show that the single relevant $d_2$-differential vanishes.  Rather than using the language of spectral sequences, we will cast the computation of this $d_2$ as a computation of a certain Toda bracket $w$.  One of the main theorems of this paper, \Cref{thm:AdamsBound}, is a lower bound on the $\HFp$-Adams filtration of $w$ for each prime $p$.  Our goals in Sections \ref{sec:ThomPushout} and \ref{sec:Toda} will be to define $w$, to reduce Theorem \ref{thm:intro-main} to the computation of $w$, and to express $w$ in as convenient a form as possible.  In particular, \Cref{lem:toda-main} will express $w$ in a form amenable to Adams filtration arguments, though we postpone any serious discussion of Adams filtration to Sections \ref{sec:Finale} and later.

\subsection{The homotopy of $\sOf$}\

The main body of this section is a computation of the homotopy of the reduced suspension spectrum $\sOf$, for $n\ge 3$.  To explain our results, it is helpful to assign names to a few elements in $\pi_*(\sOf)$ and $\pi_*(\mathbb{S})$:

\begin{dfn} \label{dfn:x}
Let $$x \in \pi_{4n-1}(\sOf)$$
denote a generator of the bottom non-zero homotopy group of $\sOf$.  Since $\sOf$ is a non-unital $\E_\infty$-ring, we may speak of the class
$$x^2 \in \pi_{8n-2}(\sOf).$$
Finally, there is a class
$$J(x) \in \pi_{4n-1}(\bS),$$
defined as the composite
$S^{4n-1} \stackrel{x}{\longrightarrow} \sOf \stackrel{J}{\longrightarrow} \bS$
\end{dfn}
The remainder of this section will consist of proofs of the following facts:
\begin{enumerate}
\item The group $\pi_{8n-2}\sOf$ is isomorphic to $\mathbb{Z}/2\mathbb{Z}$, generated by the element $x^2$ of Definition \ref{dfn:x}.
    Furthermore, the group $\pi_{8n-1} \sOf$ is isomorphic to $\pi_{8n} ko \cong \mathbb{Z}$.
\item The element $xJ(x) \in \pi_{8n-2} \sOf$, defined using the right $\pi_*(\bS)$-module structure on $\pi_*(\sOf)$, is zero.
\item
Suppose $4n-1 \le \ell \le 8n-1$.  Then the image of the map
$$\pi_{\ell}(\sOf) \stackrel{J}{\longrightarrow} \pi_{\ell}(\mathbb{S})$$
is exactly $\mathcal{J}_{\ell}$.
\end{enumerate}

The first of these facts will be proved as \Cref{cor:gen-by-x2}, the second as \Cref{lem:xJ(x)null}, and the last as \Cref{thm:imJwelldefined}.
Our key tool will be the Goodwillie tower of the identity in augmented $\mathbb{E}_\infty$-ring spectra, the basic structure of which was worked out by Nick Kuhn \cite{KuhnTAQ}.
We thank Tyler Lawson for suggesting the relevance of this tower.

\begin{dfn}
Let $X$ be a spectrum and $m \ge 1$ a natural number.  We denote by $D_m(X)$ the extended power spectrum
$$D_m(X) \coloneqq (X^{\otimes m})_{h\Sigma_m}.$$
\end{dfn}

\begin{lem} \label{lem:Kuhn}
Suppose that $R$ is an $\E_\infty$-ring spectrum, equipped with an augmentation
$$\epsilon:R \to \bS.$$
Suppose further that the fiber of $\epsilon$ is $0$-connected.  Then there is a convergent tower of $\E_{\infty}$-ring spectra:
\begin{center}
  \begin{tikzcd} [column sep=tiny]
    & & D_m(\mathrm{TAQ}(R; \bS)) \ar[d] &
    & D_2(\mathrm{TAQ}(R; \bS)) \ar[d] &
    \mathrm{TAQ}(R; \bS) \ar[d] &
    \bS \ar[d, "\simeq"]  \\
    R \ar[r] &
    \cdots \ar[r] &
    P_m(R) \ar[r] &
    \cdots \ar[r] &
    P_2(R) \ar[r] &
    P_1(R) \ar[r] &
    P_0(R)
  \end{tikzcd}
\end{center}
such that the composite map $R \to P_0(R)$ is the augmentation map $\epsilon$.
%% \begin{center}
%%   \begin{tikzcd} [row sep=small]
%%     &R \arrow{d}\\
%%     &\cdots \arrow{d}\\
%%     D_n(TAQ(R;\bS)) \arrow{r} &P_n(R) \arrow{d} \\
%%     &\cdots \arrow{d} \\
%%     D_2(TAQ(R;\bS)) \arrow{r} &P_2(R) \arrow{d} \\
%%     TAQ(R;\bS) \arrow{r} &P_1(R) \arrow{d} \\
%%     \bS \arrow{r}{\simeq} & P_0(R),
%%   \end{tikzcd}
%% \end{center}
\end{lem}

\begin{proof}
See \cite[Theorem 3.10]{KuhnTAQ}.
\end{proof}

\begin{cor} \label{cor:goodwillie}
There is a convergent tower of non-unital $\mathbb{E}_\infty$ ring spectra
\begin{center}
  \begin{tikzcd}[column sep = small]
    & & D_3(\Sigma^{-1} \tau_{\ge 4n} ko) \ar[d] &
    D_2(\Sigma^{-1} \tau_{\ge 4n} ko) \ar[d] &
    \Sigma^{-1} \tau_{\ge 4n} ko \ar[d, "\simeq"] \\
    \sOf \ar[r] &
    \cdots \ar[r] &
    Q_3 \ar[r] &
    Q_2 \ar[r] &
    Q_1 
  \end{tikzcd}
\end{center}
%% \begin{center}
%%   \begin{tikzcd} [row sep = small]
%%     &\sOf \arrow{d} \\
%%     &\cdots \arrow{d} \\
%%     D_3(\Sigma^{-1} \tau_{\ge 4n} ko) \arrow{r} & Q_3 \arrow{d} \\
%%     D_2(\Sigma^{-1} \tau_{\ge 4n} ko) \arrow{r} & Q_2 \arrow{d} \\
%%     \Sigma^{-1} \tau_{\ge 4n} ko \arrow{r}{\simeq} & Q_1.
%%   \end{tikzcd}
%% \end{center}
\end{cor}

\begin{proof}
We apply the previous lemma to $R=\spOf$ with its augmentation map $\epsilon$.  Note that, since
$$R \simeq \Sigma^{\infty}_+ \Omega^{\infty} \Sigma^{-1} \tau_{\ge 4n} ko,$$
we learn from \cite[Example 3.9]{KuhnTAQ} that $\mathrm{TAQ}(R;\bS) \simeq \Sigma^{-1} \tau_{\ge 4n} ko$.  The corollary follows by setting
$Q_i = \mathrm{fib}(P_i \to P_0).$
\end{proof}

\begin{lem}\label{lemm:d2-ko-comp}
    For $n \ge 3$, the bottom two homotopy groups of $D_2 (\Sigma^{-1} \tau_{\ge 4n} ko)$ are \[\pi_{8n-2} D_2 (\Sigma^{-1} \tau_{\ge 4n} ko) \cong \Z/2\Z\] and \[\pi_{8n-1} D_2 (\Sigma^{-1} \tau_{\ge 4n} ko) \cong 0.\]

    Moreover, the generator of $\Z/2\Z \cong \pi_{8n-2} D_2 (\Sigma^{-1} \tau_{\ge 4n} ko)$ survives in the spectral sequence associated to the tower of \Cref{cor:goodwillie} to detect $x^2 \in \pi_{8n-2} \sOf$.
\end{lem}

\begin{proof}
   There is a $4n$-connected map $\Ss^{4n-1} \to \Sigma^{-1} \tau_{\ge 4n} ko$ which induces an $(8n-1)$-connected map
	\[D_2 (\Ss^{4n-1}) \to D_2 (\Sigma^{-1} \tau_{\ge 4n} ko).\]
	Thus, there is an isomorphism 
	\[\pi_{8n-2} D_2 (\Ss^{4n-1}) \cong \pi_{8n-2} D_2 (\Sigma^{-1} \tau_{\ge 4n} ko)\]
	and a surjective map 
	\[\pi_{8n-1} D_2 (\Ss^{4n-1}) \twoheadrightarrow \pi_{8n-1} D_2 (\Sigma^{-1} \tau_{\ge 4n} ko).\]
	It therefore suffices to make the desired homotopy group computations for $D_2 (\Ss^{4n-1})$.

    There is an equivalence $D_2 (\Ss^{4n-1}) \simeq \Sigma^{4n-1}\mathbb{RP}^\infty _{4n-1}$, arising from the fact that 
		\[D_2(\Ss^{4n-1}) \simeq \Ss^{(4n-1)\rho}_{hC_2}\]
		is the Thom spectrum of the bundle $(4n-1)\mathbf{1} + (4n-1)\gamma$ over $BC_2 \simeq \mathbb{RP}^{\infty}$, and \cite[Proposition V.3.1]{BMMS} computes $\mathbb{RP}_{4n-1} ^{4n+1} \simeq \Ss^{4n-1} \cup_{2} \mathrm{e}^{4n} \cup_{\eta} \mathrm{e}^{4n+1}$. We therefore determine $\pi_{8n-2} D_2 (\Ss^{4n-1}) \cong \Z/2\Z$ and $\pi_{8n-1} D_2 (\Ss^{4n-1}) \cong 0$, as desired.

    Comparison with the Goodwillie tower of $\Sigma^{\infty} _+ \Omega^{\infty} \Ss^{4n-1}$, which recovers the Snaith splitting \cite{KuhnTAQ}, shows that the generator of $\pi_{8n-2} D_2 (\Sigma^{-1} \tau_{\ge 4n} ko)$ detects
    \[ x^2 \in \pi_{8n-2} \sOf. \qedhere\]
\end{proof}

\begin{cor} \label{cor:gen-by-x2}
For $n \ge 3$, the group $\pi_{8n-2}\sOf$ is a copy of $\mathbb{Z}/2\mathbb{Z}$, generated by the element $x^2$ of Definition \ref{dfn:x}.
    Furthermore, the group $\pi_{8n-1} \sOf$ is isomorphic to $\pi_{8n} ko$.
\end{cor}

\begin{proof}
    Note that $\tau_{\le 8n-1} D_k(\Sigma^{-1} \tau_{\ge 4n} ko)$ is trivial for $k >2$.  Thus, Corollary \ref{cor:goodwillie} and \Cref{lemm:d2-ko-comp} imply the existence of a long exact sequence
$$
\begin{tikzcd}
0 \arrow{r} & \pi_{8n-1}(\sOf) \arrow{r} \arrow[draw=none]{d}[name=Z, shape=coordinate]{} & \pi_{8n-1}(\Sigma^{-1} \tau_{\ge 4n} ko) \arrow[rounded corners,to path={ -- ([xshift=2ex]\tikztostart.east)|- (Z) [near end]\tikztonodes-| ([xshift=-2ex]\tikztotarget.west)-- (\tikztotarget)}]{dll} \\
\Z/2\Z \arrow{r}{x^2} & \pi_{8n-2}(\sOf) \arrow{r} & \pi_{8n-2}(\Sigma^{-1} \tau_{\ge 4n} ko).
\end{tikzcd}
$$

We now claim that the maps
$\pi_{k}(\sOf) \to \pi_{k}(\Sigma^{-1} \tau_{\ge 4n} ko)$ are surjective.  To see this, we note that these maps are $\pi_{k}$ of the map of spectra
$$\sOf \longrightarrow \Sigma^{-1} \tau_{\ge 4n} ko$$
that is adjoint to the identity homomorphism
$$\Of \stackrel{\simeq}{\longrightarrow} \Omega^{\infty} \Sigma^{-1} \tau_{\ge 4n} ko.$$
Indeed, it follows from the combination of \cite[Example 3.9]{KuhnTAQ} and \cite[Theorem 3.10(2)]{KuhnTAQ} that the map
$$\mathrm{hofib}(\epsilon:R \to \mathbb{S}) \to \mathrm{TAQ}(R;\mathbb{S})$$
induced by the tower of \Cref{lem:Kuhn} agrees for $R=\Sigma^{\infty}_+ \Omega^{\infty} X$ with the counit $\Sigma^{\infty}_+ \Omega^{\infty}X \to X$.
In particular, the map $\sOf \to \Sigma^{-1} \tau_{\ge 4n} ko$ admits a section after applying $\Omega^{\infty}.$ 

Identifying $\pi_{8n-2} (\Sigma^{-1} \tau_{\ge 4n} ko)$ with zero, we obtain isomorphisms
    \[\pi_{8n-1} (\sOf) \cong \pi_{8n-1} (\Sigma^{-1} \tau_{\ge 4n} ko) \cong \pi_{8n} ko\]
    and
    \[\Z/2\Z \cong \pi_{8n-2} (\sOf),\]
    with the latter sending the generator to $x^2$, as desired.
\end{proof}

\begin{cnstr}
Recall that the element $J(x) \in \pi_{4n-1} \bS$ was defined, in Definition \ref{dfn:x}, as the composite
$$S^{4n-1} \stackrel{x}{\longrightarrow} \sOf \stackrel{J}{\longrightarrow} \bS.$$
The right $\pi_*(\bS)$-module structure on $\pi_*(\sOf)$ allows us to define an element
$$xJ(x) \in \pi_{8n-2} \sOf.$$
\end{cnstr}

\begin{lem} \label{lem:xJ(x)null}
For $n \ge 3$, the element $xJ(x) \in \pi_{8n-2} \sOf$, defined using the right $\pi_*(\bS)$-module structure on $\pi_*(\sOf)$, is zero.
\end{lem}

\begin{proof}
By Corollary \ref{cor:gen-by-x2}, we know that $\pi_{8n-2} \sOf$ is isomorphic to $\Z/2\Z$, generated by the element $x^2$.
It follows that, if $xJ(x) \ne 0$, then $xJ(x) =x^2$.
In \Cref{rmk:x2-filt} we determine that the element $x^2$ has $\mathrm{H}\mathbb{F}_2$-Adams filtration $1$.
However, $xJ(x)$ has $\mathrm{H}\mathbb{F}_2$-Adams filtration at least that of $J(x)$.
Note now that $$x \in \pi_{4n-1} \sOf$$ is the suspension of an unstable class, and thus $J(x) \in \pi_{4n-1} \mathbb{S}$ is in $\J_{4n-1}$.
In particular, since $n\ge 3$, $J(x)$ has $\mathrm{H}\mathbb{F}_2$-Adams filtration larger than $1$.
\end{proof}

\subsection{The image of $J$ in $\pi_*(\mathbb{S})$}\ 

Classically, the phrase ``image of $J$'' in $\pi_{\ell}(\mathbb{S})$ refers to the image of the map
$$\pi_\ell(\mathrm{O}) \to \pi_{\ell}\GL_1(\mathbb{S}) \cong \pi_{\ell}(\mathbb{S}) \text{ for } \ell>0.$$
Recall that we use $\mathcal{J}_{\ell} \subseteq \pi_{\ell}(\mathbb{S})$ to denote this subset.

Unfortunately, we have introduced a second possible meaning of the phrase ``image of $J$,'' namely the image of the map
$$\pi_{\ell}(\sOf) \stackrel{J}{\longrightarrow} \pi_{\ell}(\mathbb{S}).$$

For general $\ell$, these two images may be different.  We prove here, however, that they are the same in our range of interest, and so no ambiguity has been introduced.
\begin{thm} \label{thm:imJwelldefined}
Suppose $n \ge 3$ and $4n-1 \le \ell \le 8n-1$.  Then the image of the map
$$\pi_{\ell}(\sOf) \stackrel{J}{\longrightarrow} \pi_{\ell}(\mathbb{S})$$
is exactly $\mathcal{J}_{\ell}$.
\end{thm}

\begin{proof}
This will automatically be true so long as every class in 
$\pi_{\ell} \sOf$ is the suspension of an unstable class in $\pi_{\ell} \Of.$
According to \Cref{cor:goodwillie}, there can be no difficulty unless $\ell=8n-1$ or $\ell=8n-2$.
The case $\ell=8n-1$ follows from \Cref{cor:gen-by-x2}.

To handle the case $\ell=8n-2$, we must check that $J(x^2)$ is an element of $\mathcal{J}_{8n-2}=0$.
Since $J$ is a non-unital ring map, $J(x^2)=J(x)^2$.  Now, it is known (by, e.g. \cite[Lemma 3]{Novikov}) that $\mathcal{J}_i \cdot \mathcal{J}_j \subseteq \mathcal{J}_{i+j}$ if $i,j>7$.  Using the hypothesis that $n \ge 3$, we have in short that $J(x)^2$ is an element of $\mathcal{J}_{8n-2} = 0$.
\end{proof}

%%% Local Variables:
%%% mode: latex
%%% TeX-master: "Main"
%%% eval: (visual-line-mode t)
%%% eval: (auto-fill-mode 0)
%%% End:
